% 香港与数字人 ... B2C MVP 短文（简体中文译本）
% 与 places.html / places-hong-kong.tex 同文。专有名词保留。无 em dash。
% 用 tectonic 或 XeLaTeX 编译。
\documentclass[11pt,a4paper]{ctexart}

\usepackage[margin=1in]{geometry}
\usepackage{setspace}
\usepackage{enumitem}
\usepackage[hidelinks]{hyperref}

\setCJKmainfont{Songti SC}
\setmainfont{Times New Roman}

\onehalfspacing
\setlength{\parskip}{0.55em}
\setlength{\parindent}{0pt}
\setlist[itemize]{leftmargin=1.4em, itemsep=0.35em, topsep=0.4em}
\Urlmuskip=0mu plus 1mu

\title{香港与数字人}
\author{Jade Zhao · 赵郎溪}
\date{2026年9月28日 \\ \small 笔记门：\texttt{zhao-langxi} · 非同行评议}

\begin{document}
\maketitle

\noindent\textit{笔记 · 香港 · 非同行评议}

看起来像人，并不是产品本身。我是印第安纳大学 Luddy 的数字人项目负责人（Digital Humans Project Lead）。
如果你是可能去跟那张脸说话的人，检验很简单：你信不信、能不能看清它是什么、你会不会再来。

\section*{你会遇到什么}

这些脸大多不是新朋友。它们坐在问询台：多语言自助机、客服屏幕、知识库上的说话层。
你提问。你要清楚的答案。你要知道自己不是在跟真人说话。
那是消费者的一天，不是销售稿，也不是白皮书。

\section*{香港目录作证据}

香港的公共目录已经把数字人列为客服与问询工具，并涵盖粤语、普通话与英语。
那是公开可见的面向个人的服务形态。不是我在香港做过的工作，也不是我在卖的产品。

数码港智慧政府创新实验室（Smart Government Innovation LAB）与香港生产力促进局的 Digital DIY 门户发布这些条目。几条公开链接：

\begin{itemize}
  \item Smart LAB · AI Digital human customer service（中国移动香港）:
    \url{https://www.smartlab.gov.hk/en/ai_solutions/a-0073}
  \item Smart LAB · Tencent Cloud Digital Human:
    \url{https://www.smartlab.gov.hk/en/ai_solutions/a-0077}
  \item HKPC Digital DIY · DYXnet Digital Human Platform:
    \url{https://ddiy.hkpc.org/en/node/1519}
  \item DYXnet 产品页（香港）:
    \url{https://www.dyxnet.com/hk/digital-human-platform/}
\end{itemize}

老实读：你看到的很多仍是视频、主播，以及服务厅问答加上一张脸。
对日常问询有用。离信任、披露、以及你会不会再来，还有距离。

\section*{这篇笔记不声称什么}

\begin{itemize}
  \item 不是说我在香港工作过。
  \item 不是产品名、价格、上线、用户数或指标。
  \item 不是说人们已经信任数字人。
  \item 不是同行评议研究。
\end{itemize}

\vspace{1.2em}
\noindent\small 独立笔记 \ldots{} 非 IU 背书研究，亦非同行评议出版物。\\
Jade Zhao · 赵郎溪 · Digital Humans Project Lead · 预计 2027 年 5 月毕业

\end{document}
